\section{Observer wave packets and spectral normalization}
\label{app:normalization}
This appendix supplies the positive/negative frequency normalization without
formally multiplying generalized eigenfunctions and delta functions.
Fix $k=|n|$, $h=k+1/2$, and $r(s)=s\tanh\pi s$. In the phase convention
of~\cite[(3.11), (4.9), (A.5)--(A.6)]{KL},
\begin{align}
 v_s(t)&:=\sqrt{\cosh t}\,\Psi_n^s(t,0)
 =N_k(s)\cosh^h t\,
 \mathbf F(h-is,h+is;k+1;z(t)),\\
 z(t)&=\frac{1-i\sinh t}{2},\\
 N_k(s)&=\sqrt{\frac{r(s)}{\pi}}\,
 e^{i\phi_{-k,s}/2}
 \frac{(1-i)(2i)^{-k}\Gamma(1/2-is)}
      {2\sqrt{2\cosh\pi s}\,\Gamma(1/2-k-is)}.
\end{align}
Here $\mathbf F={}_2F_1/\Gamma(c)$ and the half phase is chosen smoothly
and compatibly with the $n\leftrightarrow-n$ relation. Gamma recurrence
gives
\begin{equation}
 |N_k(s)|^2=\frac{r(s)2^{-2k}}{4\pi^2}
                         |\Gamma(h+is)|^2.
\end{equation}
For $t=\epsilon\tau$, $\epsilon=\pm1$, $\tau>0$, set $x=e^{-\tau}$.
The identity $-z(\epsilon\tau)=i\epsilon(1+i\epsilon x)^2/(4x)$ and
the hypergeometric connection formula~\cite[15.8.2]{DLMF} give
\begin{align}
 v_s(\epsilon\tau)&=\sum_{\sigma=\pm1}
       A_{\epsilon\sigma}(s)e^{i\sigma s\tau}
       M_{\epsilon\sigma}(x,s),\\
 A_{\epsilon\sigma}(s)&=B_\sigma(s)
          e^{-i\epsilon\pi h/2}e^{-\epsilon\sigma\pi s/2},\\
 B_\sigma(s)&=
 N_k(s)2^{h-2i\sigma s}
 \frac{\Gamma(2i\sigma s)}
      {\Gamma(h+i\sigma s)\Gamma(1/2+i\sigma s)},\\
 M_{\epsilon\sigma}(x,s)&=(1+x^2)^h
 (1+i\epsilon x)^{-2q_\sigma}
 {}_2F_1\left(q_\sigma,\tfrac12-i\sigma s;1-2i\sigma s;
       \frac{4i\epsilon x}{(1+i\epsilon x)^2}\right),\\
 q_\sigma&=h-i\sigma s.
\end{align}
The linear terms in the last expression cancel, so $M=1+x^2Q$ with
$Q$ analytic near $x=0$ on every compact positive spectral interval.
For fixed spectral and time derivative orders,
$|\partial_s^j\partial_\tau^l(M-1)|\le Ce^{-2\tau}$ for
$\tau\ge\log32$. Including the phase in the differentiated remainder
adds a factor $(1+\tau)^j$.
Thus $\int f(s)v_s(t)\dd s$ is Schwartz for $f\in\Cc((0,\infty))$.

The standard gamma moduli give
\begin{equation}
 |\Gamma(2is)|^2=\frac{\pi}{2s\sinh2\pi s},\qquad
 |\Gamma(1/2+is)|^2=\frac{\pi}{\cosh\pi s},\qquad
 |B_\sigma(s)|^2=\frac{1}{8\pi^2\cosh\pi s}.
 \label{eq:B-modulus}
\end{equation}
Now use the unitary radial variable $r_s(t)=\sqrt{2\pi}\,v_s(t)$ and
write $W_f=\int r_s f(s)\dd s$. Its Schwartz property and differential
equation imply
\begin{equation}
 e^{-iTH_k}W_f=W_{e^{-iTs^2}f}.
\end{equation}
As $T\to+\infty$, on each half-axis only the $\sigma=+$ outgoing
free packet remains. The $\sigma=-$ phase has no stationary point there,
and repeated spectral integration by parts bounds its $L^2$ norm by a
quantity tending to zero. The exponentially small Jost remainder tends
to zero on compact positions by nonstationary phase, and its tail is
uniformly $L^2$ dominated. The outgoing packet restricted to that half-axis
has the full-line Fourier norm in the limit; its omitted half-line is
again nonstationary. Plancherel therefore gives the packet Gram density
\begin{equation}
 4\pi^2\sum_{\epsilon=\pm1}
       \overline{A_{\epsilon,+}(s)}A'_{\epsilon,+}(s).
 \label{eq:packet-gram}
\end{equation}
For two positive columns, \eqref{eq:B-modulus} makes this density one.
For the negative-frequency column,
$A'_{\epsilon,+}=(-1)^nA_{-\epsilon,+}$. Its cross density is
\begin{equation}
 8\pi^2(-1)^n|B_+(s)|^2\cos(\pi h)=0.
\end{equation}
The negative column has the same unit norm. Different $n$ are already
orthogonal in $\theta$.

Finally the two asymptotic coefficient columns at $+\infty$ have determinant
\[
 -2(-1)^nB_+B_-\sinh\pi s\ne0.
\]
They span the two-dimensional
continuous solution space and are a smooth invertible change of basis
from the scattering modes on every compact positive interval. The bound
modes are orthogonal to these packets by the self-adjoint radial spectral
decomposition. This proves the full continuous normalization and
completeness needed in Section~\ref{sec:physical}; it makes no assertion
that the entire position representation lacks discrete series.
